% DERIVED from formal_layer.json — read-only, do not edit.
\begin{assumptionv}{ass:iid-subjects}[Independent identically distributed subjects]
The subject sequence \(\{(Y_i(0),Y_i(1),W_i):i\ge1\}\), comprising control potential outcomes, treatment potential outcomes, and treatment assignments, is independent and identically distributed. Equivalently, for every \(n\ge1\), the first \(n\) subject triples are jointly independent and have a common law. Zero-based indexing \(i\in\mathbb N\) corresponds to subject indexing by the relabeling \(i\leftrightarrow i+1\).
\end{assumptionv}

\begin{assumptionv}{ass:random-assignment}[Independence of treatment assignment]
For every subject \(i\ge1\), the treatment assignment \(W_i\) is independent of the potential outcomes \(Y_i(0)\) and \(Y_i(1)\):
\[
W_i\mathrel{\perp}(Y_i(0),Y_i(1)).
\]
\end{assumptionv}

\begin{assumptionv}{ass:known-assignment}[Known treatment assignment probability]
For every subject \(i\ge1\), the treatment assignment \(W_i\) has known treatment probability \(p\):
\[
P(W_i=1)=p.
\]
\end{assumptionv}

\begin{assumptionv}{ass:consistency}[Consistency of observed outcomes]
For every subject \(i\ge1\), the observed outcome \(Y_i\) is determined by the treatment assignment \(W_i\) and the treatment and control potential outcomes \(Y_i(1)\) and \(Y_i(0)\):
\[
Y_i=W_iY_i(1)+(1-W_i)Y_i(0).
\]
\end{assumptionv}

\begin{assumptionv}{ass:binary-outcomes}[Binary potential outcome support]
For every subject \(i\ge1\), the control and treatment potential outcomes \(Y_i(0)\) and \(Y_i(1)\) satisfy
\[
(Y_i(0),Y_i(1))\in\{0,1\}^2
\]
almost surely.
\end{assumptionv}

\begin{assumptionv}{ass:interior-assignment}[Interior treatment assignment probability]
The treatment assignment probability \(p\) satisfies
\[
0<p<1.
\]
\end{assumptionv}

\begin{assumptionv}{ass:interior-means}[Interior arm means]
The arm-mean parameter \(\theta\) satisfies
\[
\theta\in(0,1)^2.
\]
\end{assumptionv}

\begin{assumptionv}{ass:fixed-privacy}[Fixed positive privacy budget]
The privacy budget is a fixed real number \(\varepsilon>0\), and the privacy parameter at every sample size \(n\in\mathbb{N}\) equals \(\varepsilon\).
\end{assumptionv}

\begin{assumptionv}{ass:sequential-ldp}[Sequential local privacy]
The release rules \(Q_i\) satisfy
\[
Q_i(A\mid x,h)\le \exp(\varepsilon)Q_i(A\mid x',h)
\]
for every stage \(i\), every measurable output event \(A\in\mathcal G_i\), every pair of private records \(x,x'\in\mathcal X\), and every history tuple \(h\in\prod_{j<i}\mathcal Z_j\), including null or otherwise unattainable histories in the full history product.
\end{assumptionv}

\begin{definitionv}{def:binary-trial-model}[Binary-trial model \(\mathcal M\)]
The model is
\[
\mathcal M
=
\left\{
(P_\theta^{\otimes n})_{n\ge1}:
\theta=(\mu_0,\mu_1)^{\top}\in(0,1)^2,\;
P_\theta(X_i=x_j)=\pi_{\theta,j}
\right\},
\]
under \cref{obj:ass:interior-means,obj:ass:interior-assignment}.
Here \(\mu_0\) and \(\mu_1\) are the control-arm and treatment-arm means, \(X_i\) is subject \(i\)'s private record, \(x_j\) is an enumerated record value, and \(\pi_{\theta,j}\) is its probability under the record law \(P_\theta\). The law \(P_\theta^{\otimes n}\) is the product law of \(n\) subject records.
\end{definitionv}

\begin{definitionv}{def:sequential-channels}[Sequential private protocols $\mathfrak Q^{\mathrm{seq}}_{n,\varepsilon}$]
The class \(\mathfrak Q^{\mathrm{seq}}_{n,\varepsilon}\) consists of \(n\)-stage protocols
\[
Q^{(n)}=(Q_i)_{i=1}^n
\]
on arbitrary measurable output spaces \((\mathcal Z_i,\mathcal G_i)\), subject to \cref{obj:ass:fixed-privacy,obj:ass:sequential-ldp}, with the following properties:
\begin{itemize}
\item \textbf{(Kernels.)} Each \(Q_i(\cdot\mid x,h)\) is a Markov kernel from \(\mathcal X\times\prod_{j<i}\mathcal Z_j\) to \((\mathcal Z_i,\mathcal G_i)\).
\item \textbf{(Nonanticipation.)} Each release rule is nonanticipating.
\item \textbf{(Privacy.)} For every stage \(i\), measurable set \(A\in\mathcal G_i\), inputs \(x,x'\in\mathcal X\), and history \(h\in\prod_{j<i}\mathcal Z_j\),
\[
Q_i(A\mid x,h)\le e^\varepsilon Q_i(A\mid x',h).
\]
\end{itemize}
\end{definitionv}

\begin{definitionv}{synth_8}[Ordered staircase patterns]
For the ordered private-record alphabet
\[
\mathcal X=\{x_0,x_1,x_2,x_3\}
=\{(0,0),(0,1),(1,0),(1,1)\},
\]
let \(\mathcal S=\{S:\varnothing\ne S\subsetneq\{0,1,2,3\}\}\).
Define its fixed enumeration by
\[
\begin{aligned}
(S_1,\ldots,S_{14})=\bigl(
&\{0\},\{1\},\{0,1\},\{2\},\{0,2\},\{1,2\},\{0,1,2\},\\
&\{3\},\{0,3\},\{1,3\},\{0,1,3\},\{2,3\},\{0,2,3\},\{1,2,3\}
\bigr).
\end{aligned}
\]
Thus \(S_j\) consists of the positions of the ones in the four-bit binary expansion of \(j\), with position \(0\) the least significant bit. A pattern \(S\) corresponds to the record subset \(\{x_k:k\in S\}\). This order indexes all staircase-weight coordinates and certified active sets.
\end{definitionv}

\begin{definitionv}{synth_4}[Directional information objective]
Fix \(\theta=(\mu_0,\mu_1)^\top\in(0,1)^2\), \(p\in(0,1)\), and \(\varepsilon>0\). Write \(q=1-p\), \(d=e^\varepsilon-1\), and
\[
\pi_\theta=\bigl(q(1-\mu_0),q\mu_0,p(1-\mu_1),p\mu_1\bigr)^\top.
\]
For \(S\in\mathcal S=\{S:\varnothing\ne S\subsetneq\{0,1,2,3\}\}\), define
\[
b_S=\mathbf1+d\mathbf1_S,\qquad
h_S(\theta)=\pi_\theta^\top b_S,\qquad
g_S=d\begin{pmatrix}
q\{\mathbf1_S(1)-\mathbf1_S(0)\}\\
p\{\mathbf1_S(3)-\mathbf1_S(2)\}
\end{pmatrix}.
\]
For feasible weights
\[
\alpha\in\mathcal A_\varepsilon
=\left\{\alpha\in\mathbb R_+^{14}:
\sum_{S\in\mathcal S}\alpha_S b_S=\mathbf1\right\}
\]
and \(t\in\mathbb R\), set \(v(t)=(t,t+1)^\top\). The directional information objective is
\[
F_\theta(\alpha,t)
=\sum_{S\in\mathcal S}\alpha_S
\frac{(g_S^\top v(t))^2}{h_S(\theta)}
=v(t)^\top I_\theta(\alpha)v(t),
\qquad
I_\theta(\alpha)=
\sum_{S\in\mathcal S}\alpha_S\frac{g_Sg_S^\top}{h_S(\theta)}.
\]
Its nuisance-profiled value at \(\alpha\) is \(\min_{t\in\mathbb R}F_\theta(\alpha,t)\); the dependence on \(p\) and \(\varepsilon\) enters through \(\pi_\theta\), \(g_S\), \(b_S\), and \(\mathcal A_\varepsilon\).
\end{definitionv}

\begin{definitionv}{def:staircase-saddle}[Profiled information $J^*$ and variance $V^*$]
The optimized profiled information and its reciprocal variance are
\[
J^*(\theta,p,\varepsilon)
=
\max_{\alpha\in\mathcal A_\varepsilon}
\min_{t\in\mathbb R}
\sum_{S\in\mathcal S}
\alpha_S\frac{(g_S^{\top}v(t))^2}{h_S(\theta)},
\qquad
V^*(\theta,p,\varepsilon)
=
\frac{1}{J^*(\theta,p,\varepsilon)}.
\]
\end{definitionv}

\begin{theoremv}{thm:finite-oracle}[Finite staircase oracle]
Let \(\theta=(\mu_0,\mu_1)^\top\), let \(p\) be the treatment-assignment probability, and let \(\varepsilon\) be the privacy parameter. Suppose:
\begin{itemize}
\item \textbf{(Assignment.)} \(p\) satisfies \cref{obj:ass:interior-assignment}.
\item \textbf{(Means.)} \(\theta\) satisfies \cref{obj:ass:interior-means}.
\item \textbf{(Fixed privacy.)} The privacy sequence and \(\varepsilon\) satisfy \cref{obj:ass:fixed-privacy}.
\end{itemize}
Use the staircase objective and optimized values \(J^*(\theta,p,\varepsilon)\) and \(V^*(\theta,p,\varepsilon)=J^*(\theta,p,\varepsilon)^{-1}\) from \cref{obj:def:staircase-saddle}. For a stationary channel \(Q\), let \(I_\theta(Q)\) be the output Fisher-information matrix, obtained by integrating the products of output-density derivatives divided by the output density against the channel's dominating measure. Evaluate its contrast variance as
\[
c^\top I_\theta(Q)^+c
=
\begin{cases}
\displaystyle
\inf_{\{v:\,I_\theta(Q)v=c\}}c^\top v,
& c\in\operatorname{range}(I_\theta(Q)),\\
+\infty,
& c\notin\operatorname{range}(I_\theta(Q)).
\end{cases}
\]
For every measurable output space and every Markov channel \(Q\) on the four-record alphabet satisfying
\[
Q(A\mid x)\le e^\varepsilon Q(A\mid x')
\quad
\text{for every measurable }A\text{ and every }x,x',
\]
the oracle bound is
\[
V^*(\theta,p,\varepsilon)\le c^\top I_\theta(Q)^+c.
\]
There exists a vector \(\alpha\) of fourteen nonnegative staircase-pattern weights, with all four channel row sums equal to one, such that its staircase channel \(Q_\alpha\) satisfies the same privacy inequalities and
\[
c^\top I_\theta(Q_\alpha)^+c
=
c^\top I_\theta(\alpha)^+c
=
V^*(\theta,p,\varepsilon),
\qquad
I_\theta(\alpha)
=
\sum_{S\in\mathcal S}
\alpha_S\frac{g_Sg_S^\top}{h_S(\theta)}.
\]
Here the pattern vectors and denominators are those entering the staircase objective in \cref{obj:def:staircase-saddle}. Consequently,
\[
\inf_Q c^\top I_\theta(Q)^+c=V^*(\theta,p,\varepsilon),
\]
where the infimum ranges over all such stationary private channels and measurable output spaces. Moreover,
\[
J^*(\theta,p,\varepsilon)
=
\inf_{t\in\mathbb R}
\max_{\alpha\in\mathcal A_\varepsilon}F_\theta(\alpha,t)
=
\inf_{t\in\mathbb R}
\max_{\alpha\in\operatorname{vert}(\mathcal A_\varepsilon)}
F_\theta(\alpha,t),
\]
and there exists \(t^*\in\mathbb R\) such that
\[
\max_{\alpha\in\mathcal A_\varepsilon}
F_\theta(\alpha,t^*)
=
J^*(\theta,p,\varepsilon).
\]
\end{theoremv}

\begin{definitionv}{synth_2}[Local parameter alternative]
The local alternative \(\theta_{n,h}\) is defined coordinatewise by
\[
\theta_{n,h}(k)=\theta_0(k)+\frac{h(k)}{\sqrt n},
\qquad k\in\{0,1\},
\]
for a baseline vector \(\theta_0\in\mathbb R^2\), a perturbation vector \(h\in\mathbb R^2\), and a nonnegative integer \(n\). At \(n=0\), the quotient is defined to be zero, so \(\theta_{0,h}=\theta_0\).
\end{definitionv}

\begin{definitionv}{synth_3}[Scaled local estimation error]
The scaled local estimation error \(U_{n,h}^{\mathcal P}\) is defined by
\[
U_{n,h}^{\mathcal P}(z)
=
\sqrt n\left(
\widehat\tau_n^{\mathcal P}(z)
-
\left[
\theta_0(1)+\frac{h(1)}{\sqrt n}
-\theta_0(0)-\frac{h(0)}{\sqrt n}
\right]
\right).
\]
Here \(\mathcal P\) is a sequence of sequential \(\varepsilon\)-locally private protocols, their transcript laws factorizing through the release channels, and measurable real-valued estimators; \(\widehat\tau_n^{\mathcal P}(z)\) denotes its estimate from transcript \(z\). The parameters are \(\varepsilon\in\mathbb R\), \(\theta_0,h\in\mathbb R^2\), and \(n\in\{0,1,\ldots\}\), with \(z\) any element of the product of the \(n\) measurable release spaces. At \(n=0\), each quotient \(h(k)/\sqrt n\) is defined to be zero.
\end{definitionv}

\begin{definitionv}{synth_6}[Local alternatives and regular sequential private procedures]
Fix \(p\in(0,1)\), \(\varepsilon>0\), and \(\theta_0\in(0,1)^2\). For \(n\ge1\) and \(h\in\mathbb R^2\), set
\[
\theta_{n,h}=\theta_0+\frac{h}{\sqrt n}.
\]
For \(0<H<\infty\), the admissible local perturbations are
\[
\mathcal H_{n,H}(\theta_0)
=\left\{h\in\mathbb R^2:
\|h\|_2\le H,\ 
0<\theta_{0,k}+\frac{h_k}{\sqrt n}<1
\text{ for }k=0,1\right\}.
\]

A sequential private procedure sequence is
\[
\mathcal P=((Q^{(n)},\widehat\tau_n))_{n\ge1},
\]
where \(Q^{(n)}\in\mathfrak Q^{\mathrm{seq}}_{n,\varepsilon}\) for every \(n\), and \(\widehat\tau_n\) is \(\sigma(Z_1,\ldots,Z_n)\)-measurable. With \(c=(-1,1)^\top\), its scaled local error is
\[
U_{n,h}^{\mathcal P}
=\sqrt n\{\widehat\tau_n-c^\top\theta_{n,h}\}.
\]
We say that \(\mathcal P\) is regular at \(\theta_0\) when there exists a Borel probability law \(L_{\mathcal P}\) on \(\mathbb R\) satisfying
\[
\int u^2\,dL_{\mathcal P}(u)<\infty
\]
such that, for every fixed \(h\in\mathbb R^2\),
\[
U_{n,h}^{\mathcal P}\rightsquigarrow L_{\mathcal P}
\]
under the experiment at \(\theta_{n,h}\), with the same law \(L_{\mathcal P}\) for every fixed \(h\), and the local squared errors are asymptotically uniformly integrable: for every \(0<H<\infty\),
\[
\lim_{M\to\infty}\limsup_{n\to\infty}
\sup_{h\in\mathcal H_{n,H}(\theta_0)}
\mathbb E_{\theta_{n,h}}^{\mathcal P}
\left[
(U_{n,h}^{\mathcal P})^2
\mathbf1\{|U_{n,h}^{\mathcal P}|>M\}
\right]=0.
\]
For fixed \(h\), the local alternative is interior for all sufficiently large \(n\). The class \(\mathfrak P^{\mathrm{reg}}_{\varepsilon}(\theta_0)\) consists of all complete procedure sequences satisfying these requirements.
\end{definitionv}

\begin{assumptionv}{ass:pilot-diverges}[Diverging pilot size]
The pilot sample size \(m_n\) satisfies
\[
m_n\to\infty.
\]
\end{assumptionv}

\begin{assumptionv}{ass:pilot-sublinear}[Sublinear pilot sample size]
The pilot sample sizes \(m_n\) are nonnegative integers satisfying
\[
\frac{m_n}{n}\longrightarrow 0
\quad\text{as }n\to\infty,
\qquad
1\le m_n<n
\quad\text{for every integer }n\ge2.
\]
\end{assumptionv}

\begin{definitionv}{synth_10}[Four-category randomized-response pilot]
For \(\varepsilon>0\) and the ordered private-record categories
\[
x_0=(0,0),\qquad x_1=(0,1),\qquad
x_2=(1,0),\qquad x_3=(1,1),
\]
four-category \(\varepsilon\)-randomized response is the release kernel on \(\{0,1,2,3\}\) defined by
\[
Q_{\mathrm{RR}}(k\mid x_j)
=
\begin{cases}
\dfrac{\exp(\varepsilon)}{\exp(\varepsilon)+3},&k=j,\\[6pt]
\dfrac{1}{\exp(\varepsilon)+3},&k\ne j,
\end{cases}
\qquad j,k\in\{0,1,2,3\}.
\]
The pilot releases \(R_1,\ldots,R_{m_n}\) are generated by applying this kernel to the first \(m_n\) private records with independent randomization. Consequently, for each reported category \(k\),
\[
\Pr_\theta(R_i=k)
=\frac{1+(\exp(\varepsilon)-1)\pi_{\theta,k}}
       {\exp(\varepsilon)+3}.
\]
\end{definitionv}

\begin{algorithmv}{def:pilot-estimator}[Private pilot estimator]
The private pilot estimator is the sequence of sequential \(\varepsilon\)-locally private protocols and measurable estimators defined below, for an assignment probability \(p\) satisfying \cref{obj:ass:interior-assignment}, a privacy parameter \(\varepsilon>0\), and an arbitrary pilot-size schedule \(m:\mathbb N\to\mathbb N\), with \(m_n=m(n)\).
\begin{enumerate}
\item Fix a measurable selector defined on all real parameter pairs,
\[
\vartheta\longmapsto
(\alpha_{\vartheta},t_{\vartheta},J_{\vartheta}),
\]
whose weight component has fourteen real coordinates. For every interior parameter \(\vartheta\), its values satisfy
\[
\begin{aligned}
\alpha_{\vartheta}&\in\mathcal A_\varepsilon,\\
t_{\vartheta}&\in\arg\min_{t\in\mathbb R}
F_{\vartheta}(\alpha_{\vartheta},t),\\
J_{\vartheta}
&=F_{\vartheta}(\alpha_{\vartheta},t_{\vartheta})
=J^*(\vartheta,p,\varepsilon),\\
F_{\vartheta}(\alpha,t_{\vartheta})
&\le J_{\vartheta}
\qquad\text{for every }\alpha\in\mathcal A_\varepsilon.
\end{aligned}
\]
The staircase quantities are as in \cref{obj:def:staircase-saddle}.

\item At sample size \(n\), release each available record with index \(i\le m_n\) through four-category \(\varepsilon\)-randomized response, obtaining pilot releases \(R_i\in\{0,1,2,3\}\), with the category indexing of \cref{obj:synth_10}. Define
\[
\widehat u_k
=\frac{1}{m_n}
\sum_{i=1}^{\min\{m_n,n\}}\mathbf 1\{R_i=k\},
\qquad
\widehat\pi_k
=\frac{(\exp(\varepsilon)+3)\widehat u_k-1}
{\exp(\varepsilon)-1},
\qquad k\in\{0,1,2,3\}.
\]
Reciprocals of zero are interpreted as zero, and empty sums as zero.

\item Define the clipping margin and pilot parameter estimate by
\[
\begin{aligned}
\delta_n&=(m_n+2)^{-1},\\
\widetilde\theta_0
&=\max\left\{\delta_n,
\min\left\{1-\delta_n,\frac{\widehat\pi_1}{1-p}\right\}\right\},\\
\widetilde\theta_1
&=\max\left\{\delta_n,
\min\left\{1-\delta_n,\frac{\widehat\pi_3}{p}\right\}\right\},\\
\widetilde\theta
&=(\widetilde\theta_0,\widetilde\theta_1)^{\top}.
\end{aligned}
\]
Here \(\widehat\pi_1\) and \(\widehat\pi_3\) use the zero-based control-success and treatment-success categories, respectively. Evaluate the selector at this estimate:
\[
(\widetilde\alpha,\widetilde t,\widetilde J)
=
(\alpha_{\widetilde\theta},
t_{\widetilde\theta},
J_{\widetilde\theta}).
\]

\item Release each remaining record through the selected staircase channel:
\[
S_i\sim Q_{\widetilde\alpha}(\,\cdot\mid X_i),
\qquad m_n<i\le n,
\]
conditionally on the pilot releases and the remaining records. The transcript law is the law generated by the successive release kernels. Output
\[
\widehat\tau^*
=
c^{\top}\widetilde\theta
+\bigl(\max\{n-m_n,0\}\bigr)^{-1}
\sum_{\substack{1\le i\le n\\m_n<i}}
\frac{g_{S_i}^{\top}v(\widetilde t)}
{\widetilde J\,h_{S_i}(\widetilde\theta)}.
\]
\end{enumerate}
\end{algorithmv}

\begin{theoremv}{thm:sequential-minimax}[Sequential local minimax efficiency]
Fix the treatment-assignment probability \(p\), the privacy level \(\varepsilon\), and the arm-mean parameter \(\theta_0\). Suppose:
\begin{itemize}
\item \textbf{(Interior assignment.)} \(p\) satisfies \cref{obj:ass:interior-assignment}.
\item \textbf{(Interior means.)} \(\theta_0\) satisfies \cref{obj:ass:interior-means}.
\item \textbf{(Fixed privacy.)} The privacy sequence and \(\varepsilon\) satisfy \cref{obj:ass:fixed-privacy}.
\item \textbf{(Strong-saddle selection.)} Fix any selector \(\vartheta\mapsto(\alpha_{\vartheta},t_{\vartheta},J_{\vartheta})\) satisfying the strong-saddle selector conditions in \cref{obj:def:pilot-estimator}.
\end{itemize}
Write \(V^*(\theta_0,p,\varepsilon)=J^*(\theta_0,p,\varepsilon)^{-1}\) for the reciprocal optimized profiled information in \cref{obj:def:staircase-saddle}. For an arm-mean parameter \(\theta=(\mu_0,\mu_1)^\top\), the treatment-effect target is
\[
\tau(\theta)=\mu_1-\mu_0.
\]
For \(n\ge1\) and \(H>0\), define the local alternatives and admissible perturbations by
\[
\theta_{n,h}=\theta_0+\frac{h}{\sqrt n},
\qquad
\mathcal H_{n,H}(\theta_0)
=
\left\{
h\in\mathbb R^2:
\|h\|_2\le H,\quad
\theta_0+\frac{h}{\sqrt n}\in(0,1)^2
\right\}.
\]
For a procedure sequence \(\mathcal P\), let its scaled local estimation error be
\[
U_{n,h}^{\mathcal P}
=
\sqrt n\left(
\widehat\tau_n^{\mathcal P}-\tau(\theta_{n,h})
\right),
\]
where \(\widehat\tau_n^{\mathcal P}\) is its transcript-measurable estimator.

For every sequence \(\mathcal P\) of sequential \(\varepsilon\)-locally private protocols and measurable estimators, with arbitrary measurable output spaces indexed by sample size and release stage and with transcript laws factorizing according to the history-dependent kernels in \cref{obj:def:sequential-channels},
\[
V^*(\theta_0,p,\varepsilon)
\le
\sup_{H>0}\;\liminf_{n\to\infty}\;
\sup_{h\in\mathcal H_{n,H}(\theta_0)}
\mathbb E_{\theta_{n,h}}^{\mathcal P}
\!\left[(U_{n,h}^{\mathcal P})^2\right].
\]

Let \(\mathfrak P^{\mathrm{reg}}_{\varepsilon}(\theta_0)\) denote the class of these procedure sequences whose scaled local errors have a common local weak limit law \(L_{\mathcal P}\) with finite second moment,
\[
\int_{\mathbb R}u^2\,L_{\mathcal P}(du)<\infty,
\]
and whose squared-error tails satisfy, for every \(H>0\),
\[
\lim_{M\to\infty}\;\limsup_{n\to\infty}\;
\sup_{h\in\mathcal H_{n,H}(\theta_0)}
\mathbb E_{\theta_{n,h}}^{\mathcal P}
\!\left[
(U_{n,h}^{\mathcal P})^2
\mathbf 1\{|U_{n,h}^{\mathcal P}|>M\}
\right]
=0.
\]

There exists a pilot schedule, specifically
\[
m_n=\lfloor\sqrt n\rfloor,
\]
satisfying \cref{obj:ass:pilot-diverges,obj:ass:pilot-sublinear}, for which the procedure in \cref{obj:def:pilot-estimator}, using the fixed selector above, belongs to \(\mathfrak P^{\mathrm{reg}}_{\varepsilon}(\theta_0)\) and attains equality in the lower bound. Moreover,
\[
\inf_{\mathcal P\in\mathfrak P^{\mathrm{reg}}_{\varepsilon}(\theta_0)}
\;\sup_{H>0}\;\liminf_{n\to\infty}\;
\sup_{h\in\mathcal H_{n,H}(\theta_0)}
\mathbb E_{\theta_{n,h}}^{\mathcal P}
\!\left[(U_{n,h}^{\mathcal P})^2\right]
=
V^*(\theta_0,p,\varepsilon).
\]
All risks and the infimum are interpreted in the extended nonnegative reals, with the real variance constant embedded therein.
\end{theoremv}

\begin{definitionv}{synth_9}[Pilot-selected affine correction]
Let \(p\in(0,1)\), \(\varepsilon>0\), and let \(\widetilde\theta\in(0,1)^2\) be the private pilot estimate. Apply the measurable strong-saddle selector of \cref{obj:def:pilot-estimator} to obtain
\[
(\widetilde\alpha,\widetilde t,\widetilde J)
=(\alpha_{\widetilde\theta},t_{\widetilde\theta},J_{\widetilde\theta}),
\qquad
\widetilde J
=F_{\widetilde\theta}(\widetilde\alpha,\widetilde t)
=J^*(\widetilde\theta,p,\varepsilon)>0.
\]
For each \(S\in\mathcal S\), define the selected affine correction by
\[
\widetilde\phi(S)
=\frac{g_S^{\top}v(\widetilde t)}
       {\widetilde J\,h_S(\widetilde\theta)}.
\]
Here \(v(t)=(t,t+1)^{\top}\), \(h_S(\vartheta)=\pi_\vartheta^{\top}b_S\), and \(g_S=\nabla_\vartheta h_S(\vartheta)\). The selected weights determine the main-sample channel
\[
Q_{\widetilde\alpha}(S\mid x_j)=\widetilde\alpha_S b_S(j)
\]
and the information value
\[
\widetilde J
=\sum_{T\in\mathcal S}
\widetilde\alpha_T
\frac{(g_T^{\top}v(\widetilde t))^2}
     {h_T(\widetilde\theta)}.
\]
For a transcript entry that is a pilot release, set its correction to zero; for a main-sample pattern release \(S\), use \(\widetilde\phi(S)\).
\end{definitionv}

\begin{definitionv}{synth_1}[Empirical variance of selected corrections]
The variance estimator \(\widehat V^*\) is defined for any \(p,\varepsilon\in\mathbb R\), pilot-size schedule \(m:\mathbb N\to\mathbb N\), sample size \(n\in\mathbb N\), transcript \(z\) of \(n\) pilot-or-pattern releases, and selector mapping each real parameter pair to a staircase-weight vector and two real numbers. Write \(m_n=m(n)\), and let \(\widetilde\theta\) be the pilot parameter estimate computed from \(p,\varepsilon,m,z\). For each transcript entry, \(\widetilde\phi(z_i)\) denotes zero when the entry is a pilot release and the selected affine correction \(\widetilde\phi(S_i)\), evaluated at \(\widetilde\theta,p,\varepsilon\) using the supplied selector, when the entry is a pattern release. Then
\[
\widehat V^*
=
\frac{1}{\max\{n-m_n,0\}}
\sum_{\substack{0\le i<n\\m_n\le i}}
\left(
\widetilde\phi(z_i)
-
\frac{1}{\max\{n-m_n,0\}}
\sum_{\substack{0\le j<n\\m_n\le j}}
\widetilde\phi(z_j)
\right)^2.
\]
Here reciprocals are interpreted as zero at zero, so the definition also gives \(\widehat V^*=0\) when \(m_n\ge n\).
\end{definitionv}

\begin{theoremv}{thm:pilot-attainment}[Private pilot attainment]
Let \(\theta=(\mu_0,\mu_1)^{\top}\) be the arm-mean parameter, \(p\) the treatment-assignment probability, and \(\varepsilon\) the fixed privacy level. Suppose:
\begin{itemize}
\item \textbf{(Assignment.)} \(p\) satisfies \cref{obj:ass:interior-assignment}.
\item \textbf{(Means.)} \(\theta\) satisfies \cref{obj:ass:interior-means}.
\item \textbf{(Privacy.)} The privacy sequence and \(\varepsilon\) satisfy \cref{obj:ass:fixed-privacy}.
\item \textbf{(Selection.)} The selector in \cref{obj:def:pilot-estimator} satisfies the saddle-selection conditions there and, at every interior parameter \(\vartheta\), its selected weights maximize the objective at its selected profiling coordinate:
\[
F_{\vartheta}(\alpha,t_{\vartheta})
\le F_{\vartheta}(\alpha_{\vartheta},t_{\vartheta})
\qquad
\text{for every }\alpha\in\mathcal A_{\varepsilon}.
\]
\item \textbf{(Pilot schedule.)} The pilot sizes \(m_n\) satisfy \cref{obj:ass:pilot-diverges,obj:ass:pilot-sublinear}.
\item \textbf{(Coverage level.)} The Wald error probability \(\gamma\) satisfies \(0<\gamma<1\).
\end{itemize}

Consider the complete procedure sequence \(\mathcal P\) constructed in \cref{obj:def:pilot-estimator}, with estimator \(\widehat\tau^*\), and write \(\tau=\mu_1-\mu_0\) for the treatment-effect contrast. For every sample size, its channels satisfy sequential \(\varepsilon\)-local differential privacy as specified in \cref{obj:ass:sequential-ldp}.

For \(n\ge2\), let \(H_{m_n}=(R_1,\ldots,R_{m_n})\) be the private pilot history. For every bounded measurable function \(f\) of the transcript whose value depends only on \(H_{m_n}\), the product \((\widehat\tau^*-\tau)f\) is integrable and
\[
\mathbb E_{\theta}\!\left[(\widehat\tau^*-\tau)f\right]=0.
\]
Under every fixed local perturbation \(h\), the scaled local error \(U_{n,h}^{\mathcal P}\), centered at the contrast of the local parameter \(\theta_{n,h}\), satisfies
\[
U_{n,h}^{\mathcal P}
=
\sqrt n\left(
\widehat\tau^*
-
\bigl((\theta_{n,h})_1-(\theta_{n,h})_0\bigr)
\right)
\rightsquigarrow
N\!\left(0,V^*(\theta,p,\varepsilon)\right)
\quad\text{under the transcript law at }\theta_{n,h}.
\]
Here \(V^*(\theta,p,\varepsilon)\) is the optimized contrast variance in \cref{obj:def:staircase-saddle}. In particular,
\[
\sqrt n(\widehat\tau^*-\tau)
\rightsquigarrow N\!\left(0,V^*(\theta,p,\varepsilon)\right).
\]
The variance estimator \(\widehat V^*\) satisfies
\[
\Pr_{\theta}\!\left(
\left|\widehat V^*-V^*(\theta,p,\varepsilon)\right|>\delta
\right)\longrightarrow0
\qquad\text{for every }\delta>0.
\]

The procedure is regular at \(\theta\), with common local limit law
\[
L_{\mathcal P}=N\!\left(0,V^*(\theta,p,\varepsilon)\right).
\]
In addition to the local convergence above, its squared scaled errors satisfy asymptotic uniform integrability over every bounded admissible local-perturbation set: writing \(\mathcal H_{n,H}(\theta)\) for that set at radius \(H\), for every \(H>0\),
\[
\lim_{\substack{M\to\infty\\ M\in\mathbb N}}
\limsup_{n\to\infty}
\sup_{h\in\mathcal H_{n,H}(\theta)}
\mathbb E_{\theta_{n,h}}\!\left[
\bigl(U_{n,h}^{\mathcal P}\bigr)^2
\mathbf 1\!\left\{\left|U_{n,h}^{\mathcal P}\right|>M\right\}
\right]
=0.
\]

Finally, let \(\Phi\) be the standard normal distribution function and define its critical value by
\[
z_{1-\gamma/2}
=
\inf\left\{z\in\mathbb R:1-\gamma/2\le\Phi(z)\right\}.
\]
Then
\[
\Pr_{\theta}\!\left(
\left|\widehat\tau^*-\tau\right|
\le z_{1-\gamma/2}\sqrt{\widehat V^*/n}
\right)\longrightarrow1-\gamma.
\]
\end{theoremv}

\begin{theoremv}{thm:five-output-upper-bound}[Five-output optimal staircase channel]
Let \(\theta=(\mu_0,\mu_1)^\top\), \(p\), and \(\varepsilon\) satisfy:
\begin{itemize}
\item \textbf{(Interior means.)} \(0<\mu_0<1\) and \(0<\mu_1<1\).
\item \textbf{(Assignment probability.)} \(0<p<1\).
\item \textbf{(Privacy parameter.)} \(\varepsilon>0\).
\end{itemize}
With the feasible staircase weights and profiled objective of \cref{obj:def:staircase-saddle}, there exists \(\alpha\in\mathcal A_\varepsilon\) such that
\[
J^*(\theta,p,\varepsilon)
=
\inf_{t\in\mathbb R}F_\theta(\alpha,t),
\qquad
\bigl|\{S\in\mathcal S:\alpha_S\ne0\}\bigr|\le5,
\qquad
\kappa(Q_\alpha)\le5.
\]
Here the active-output cardinality of the associated fourteen-symbol staircase channel is
\[
\kappa(Q_\alpha)
=
\bigl|\{S\in\mathcal S:(P_\theta Q_\alpha)(\{S\})\ne0\}\bigr|.
\]
\end{theoremv}

\begin{definitionv}{synth_5}[Staircase support certificates]
Fix \(\theta\in(0,1)^2\), \(p\in(0,1)\), and \(\varepsilon>0\), and use the staircase rays, masses, and derivative vectors defined above. For a proposed active set \(A\subseteq\mathcal S\), a support certificate consists of weights \(\alpha\in\mathbb R_+^{14}\), a profiling coordinate \(t\in\mathbb R\), and a dual vector \(\eta\in\mathbb R^4\), with \(\alpha_S>0\) for \(S\in A\) and \(\alpha_S=0\) for \(S\notin A\). Write
\[
f_S(t)=\frac{(g_S^\top v(t))^2}{h_S(\theta)}.
\]
We say that primal normalization holds when
\[
B\alpha=\mathbf1,
\qquad\text{equivalently}\qquad
\sum_{S\in\mathcal S}\alpha_S b_S=\mathbf1,
\]
where \(B=(b_{S_1},\ldots,b_{S_{14}})\).

We say that profiling stationarity holds at \((\alpha,t)\) when
\[
\frac{\partial}{\partial t}F_\theta(\alpha,t)
=\sum_{S\in\mathcal S}\alpha_S f_S'(t)
=2\sum_{S\in\mathcal S}\alpha_S
\frac{(g_S^\top v(t))(g_S^\top(1,1)^\top)}{h_S(\theta)}
=0.
\]
For feasible \(\alpha\), the function \(F_\theta(\alpha,\cdot)\) is a differentiable convex quadratic, so this condition is equivalent to \(t\in\arg\min_{u\in\mathbb R}F_\theta(\alpha,u)\).

The dual vector \(\eta\) corresponds to the normalization constraint in the linear program
\[
\max_{\alpha\ge0:\,B\alpha=\mathbf1}
\sum_{S\in\mathcal S}\alpha_S f_S(t);
\]
its dual feasibility conditions are \(b_S^\top\eta\ge f_S(t)\) for every \(S\in\mathcal S\). We say that active dual equalities hold when
\[
b_S^\top\eta=f_S(t)\qquad(S\in A),
\]
and that strict inactive dual inequalities hold when
\[
b_S^\top\eta>f_S(t)\qquad(S\in\mathcal S\setminus A).
\]
The active sets for the three-output and five-output certificates are respectively
\[
A=\{S_1,S_6,S_8\},
\qquad
A=\{S_3,S_4,S_6,S_8,S_9\},
\]
using the fixed pattern enumeration. The five-output certificate additionally requires pairwise distinct values of
\[
\rho_S(\theta,t)=\frac{g_S^\top v(t)}{h_S(\theta)}
\qquad(S\in A).
\]
\end{definitionv}

\begin{definitionv}{def:support-regions}[Support regions $\mathcal R_3$ and $\mathcal R_5$]
The region \(\mathcal R_3\) is the parameter set for which, using the fixed pattern enumeration \(S_1,\ldots,S_{14}\), the following conditions hold with active set \(\{S_1,S_6,S_8\}\):
\begin{itemize}
\item \textbf{(Normalization.)} Primal normalization holds.
\item \textbf{(Stationarity.)} Profiling stationarity holds.
\item \textbf{(Active patterns.)} The active dual equalities hold.
\item \textbf{(Inactive patterns.)} The inactive dual inequalities are strict.
\end{itemize}
The region \(\mathcal R_5\) is the analogous parameter set with active set \(\{S_3,S_4,S_6,S_8,S_9\}\), additionally satisfying
\[
\rho_S(\theta,t)\ne\rho_T(\theta,t)
\qquad
\text{for all distinct }S,T\in\{S_3,S_4,S_6,S_8,S_9\}.
\]
\end{definitionv}

\begin{theoremv}{thm:five-output-region}[The five-output region]
The five-pattern certificate and score-separation region \(\mathcal R_5\), associated with the support \(\{S_3,S_4,S_6,S_8,S_9\}\), is open in \(\mathbb R^4\) and nonempty. Write \(J^*(\theta,p,\varepsilon)\) and \(V^*(\theta,p,\varepsilon)\) for the optimized profiled information and its reciprocal from \cref{obj:def:staircase-saddle}.

For every \(0<p<1/2\),
\[
(p,1/2,1/2,\log 3)\in\mathcal R_5.
\]
At \(\theta=(1/2,1/2)^\top\) and \(\varepsilon=\log 3\), there is a unique feasible staircase-weight vector \(\alpha\) attaining the optimized profiled information:
\[
J^*(\theta,p,\log 3)
=
\inf_{t\in\mathbb R}F_\theta(\alpha,t),
\qquad
\{S\in\mathcal S:\alpha_S\ne0\}
=
\{S_3,S_4,S_6,S_8,S_9\}.
\]
Here feasibility means that the weights are nonnegative and their weighted staircase rays sum to one at each of the four input records. Uniqueness means that every feasible weight vector \(\beta\) satisfying
\[
J^*(\theta,p,\log 3)
=
\inf_{t\in\mathbb R}F_\theta(\beta,t)
\]
equals \(\alpha\).

For every \(1/2<p<1\),
\[
(1-p,1/2,1/2,\log 3)\in\mathcal R_5.
\]
At \(\theta=(1/2,1/2)^\top\) and the original assignment probability \(p\), there is likewise a unique feasible staircase-weight vector \(\alpha\) satisfying
\[
J^*(\theta,p,\log 3)
=
\inf_{t\in\mathbb R}F_\theta(\alpha,t).
\]
Its support is the image of \(\{S_3,S_4,S_6,S_8,S_9\}\) under exchange of the treatment labels.

In each of these two cases, every stationary \((\log 3)\)-LDP channel \(Q\) from the four-record alphabet to an arbitrary measurable output space that attains contrast variance \(V^*(\theta,p,\log 3)\) satisfies \(\kappa(Q)\ge5\). More generally, for every
\[
(p,\mu_0,\mu_1,\varepsilon)\in\mathcal R_5,
\qquad
\theta=(\mu_0,\mu_1)^\top,
\]
every stationary \(\varepsilon\)-LDP channel \(Q\) attaining contrast variance \(V^*(\theta,p,\varepsilon)\) satisfies
\[
\kappa(Q)\ge5.
\]
A stationary \(\varepsilon\)-LDP channel is a Markov kernel satisfying
\[
Q(x,A)\le e^\varepsilon Q(x',A)
\]
for every measurable output event \(A\) and every pair of input records \(x,x'\). Attainment means that the extended nonnegative contrast variance of the output Fisher-information matrix \(I_\theta(Q)\) equals the nonnegative part of \(V^*(\theta,p,\varepsilon)\). Explicitly, this contrast variance is
\[
\begin{cases}
\displaystyle
\max\!\left\{
\inf_{v:\,I_\theta(Q)v=c}c^\top v,\;0
\right\},
& c\in\operatorname{range}(I_\theta(Q)),\\[4pt]
+\infty,
& c\notin\operatorname{range}(I_\theta(Q)),
\end{cases}
\]
where \(c\) is the treatment-effect contrast vector. The output-cardinality convention is
\[
\kappa(Q)=
\begin{cases}
\#\{z\in\mathcal Z:(P_\theta Q)(\{z\})\ne0\},
& \mathcal Z\ \text{finite},\\
+\infty,
& \mathcal Z\ \text{infinite},
\end{cases}
\]
where \(P_\theta Q\) is the channel's output law.
\end{theoremv}

\begin{theoremv}{thm:three-output-region}[Three-output region and optimizer]
The three-pattern certificate region \(\mathcal R_3\) is open in \(\mathbb R^4\) and contains
\[
(p,\mu_0,\mu_1,\varepsilon)
=\left(\frac{3}{10},\frac{1}{10},\frac{1}{10},\log 3\right).
\]
At this point, with \(\theta=(1/10,1/10)^\top\), there exists a feasible staircase-weight vector \(\alpha\), in the optimization of \cref{obj:def:staircase-saddle}, whose support and weights satisfy
\[
\{S\in\mathcal S:\alpha_S\ne0\}=\{S_1,S_6,S_8\},
\qquad
\alpha_{S_1}=\alpha_{S_6}=\alpha_{S_8}=\frac15.
\]
This vector attains the optimized profiled information:
\[
J^*\left(\theta,\frac{3}{10},\log 3\right)
=\inf_{t\in\mathbb R}F_\theta(\alpha,t),
\]
where the profiled objective is evaluated at \(p=3/10\) and \(\varepsilon=\log 3\). Moreover, every feasible staircase-weight vector \(\beta\) at \(\varepsilon=\log 3\) satisfying
\[
J^*\left(\theta,\frac{3}{10},\log 3\right)
=\inf_{t\in\mathbb R}F_\theta(\beta,t)
\]
equals \(\alpha\). The optimized information and its reciprocal variance, as defined in \cref{obj:def:staircase-saddle}, are
\[
J^*\left(\theta,\frac{3}{10},\log 3\right)=\frac{441}{3994},
\qquad
V^*\left(\theta,\frac{3}{10},\log 3\right)=\frac{3994}{441}.
\]
\end{theoremv}

\begin{theoremv}{thm:support-transition}[Local support and cardinality transition]
Fix the treatment-assignment probability \(p=3/10\) and privacy parameter \(\varepsilon=\log 3\). There exist disjoint open sets \(U_5,U_3\subseteq\mathbb R^2\) containing \((1/2,1/2)\) and \((1/10,1/10)\), respectively, with the following properties:
\begin{itemize}
\item For every \(\theta=(\mu_0,\mu_1)^\top\in U_5\), there exists a unique feasible staircase-weight vector \(\alpha\in\mathcal A_\varepsilon\) satisfying
\[
J^*(\theta,p,\varepsilon)
=
\inf_{t\in\mathbb R}F_\theta(\alpha,t).
\]
Its active support is \(\{S_3,S_4,S_6,S_8,S_9\}\). Moreover, a feasible staircase channel attains contrast variance \(\max\{V^*(\theta,p,\varepsilon),0\}\) with exactly five active outputs, and every stationary \(\varepsilon\)-LDP channel on any measurable output space attaining this contrast variance has at least five active outputs.
\item For every \(\theta=(\mu_0,\mu_1)^\top\in U_3\), there exists a unique feasible staircase-weight vector \(\alpha\in\mathcal A_\varepsilon\) satisfying
\[
J^*(\theta,p,\varepsilon)
=
\inf_{t\in\mathbb R}F_\theta(\alpha,t).
\]
Its active support is \(\{S_1,S_6,S_8\}\). Moreover, a feasible staircase channel attains contrast variance \(\max\{V^*(\theta,p,\varepsilon),0\}\) with exactly three active outputs, and every stationary \(\varepsilon\)-LDP channel on any measurable output space attaining this contrast variance has at least three active outputs.
\end{itemize}
In each case, uniqueness is among all feasible staircase-weight vectors satisfying the displayed equality.
\end{theoremv}

\begin{theoremv}{thm:balanced-reduction}[Balanced randomized response reduction]
Let \(\theta=(\mu_0,\mu_1)^\top\) be the arm-mean parameter and \(\tau=\mu_1-\mu_0\) the treatment-effect contrast. Suppose:
\begin{itemize}
\item \textbf{(Interior means.)} \(\theta\) satisfies \cref{obj:ass:interior-means}.
\item \textbf{(Balance.)} \(p=1/2\) and \(\mu_0+\mu_1=1\).
\item \textbf{(Privacy parameter.)} \(\varepsilon>0\).
\end{itemize}
Let \(Q\) be binary randomized response applied to \((2W_i-1)(2Y_i-1)\): it releases the input sign with probability
\[
\frac{e^\varepsilon}{e^\varepsilon+1}
\]
and the opposite sign with probability
\[
\frac{1}{e^\varepsilon+1}.
\]
Write \(V^*(\theta,p,\varepsilon)=J^*(\theta,p,\varepsilon)^{-1}\) for the reciprocal optimized nuisance-profiled information. Then
\[
V^*(\theta,1/2,\varepsilon)
=
\left(\frac{\cosh(\varepsilon/2)}{\sinh(\varepsilon/2)}\right)^2-\tau^2.
\]
The channel \(Q\) is a stationary Markov channel satisfying, for every measurable output event and every pair of input records,
\[
Q(x_j,A)\le e^\varepsilon Q(x_{j'},A).
\]
Moreover, let \(I_\theta(Q)\) be the Fisher-information matrix of its output experiment at \(p=1/2\), and let \(c=(-1,1)^\top\). With contrast variance defined as \(c^\top I_\theta(Q)^+c\) when \(c\in\operatorname{range}(I_\theta(Q))\), and as \(+\infty\) otherwise, this channel attains
\[
c^\top I_\theta(Q)^+c=V^*(\theta,1/2,\varepsilon),
\qquad
c\in\operatorname{range}(I_\theta(Q)).
\]
\end{theoremv}

\begin{definitionv}{def:oa-release}[Custom-Laplace release $\widetilde A_i$]
The custom-Laplace release and its sample-mean estimator are
\[
\widetilde A_i
=
\frac{W_iY_i}{p}
-
\frac{(1-W_i)Y_i}{q}
+
L_i,
\qquad
\widehat\tau_{\mathrm{OA}}
=
\frac{1}{n}\sum_{i=1}^{n}\widetilde A_i.
\]
The noise variables \(L_i\) are independent and identically distributed, independent of all subject records, with
\[
L_i\sim\operatorname{Lap}(0,\Delta_A/\varepsilon),
\qquad
\Delta_A=p^{-1}+q^{-1}=\frac{1}{pq}.
\]
\end{definitionv}

\begin{theoremv}{thm:laplace-comparison}[Laplace variance and interval comparison]
Let \(\theta=(\mu_0,\mu_1)^\top\) be the pair of arm means, \(p\) the treatment-assignment probability, \(\varepsilon\) the privacy parameter, and \(\gamma\) the Wald error probability. Suppose the following conditions hold, under the usual measurability conditions.
\begin{itemize}
\item \textbf{(Binary randomized trial.)} The subjects \(\{(Y_i(0),Y_i(1),W_i):i\ge1\}\) are independent and identically distributed. The potential outcomes are binary, treatment assignment is independent of the potential outcomes, and
\[
\Pr(W_i=1)=p,\qquad
Y_i=(1-W_i)Y_i(0)+W_iY_i(1),\qquad
\mu_0=\mathbb E[Y_i(0)],\qquad
\mu_1=\mathbb E[Y_i(1)],
\]
with \(0<p<1\) and \(0<\mu_0,\mu_1<1\). Set \(q=1-p\) and \(\tau=\mu_1-\mu_0\).
\item \textbf{(Public experiment and representation.)} The specified public experiment is \(P_\theta^{\otimes n}\) for every positive sample size \(n\), and the private records \(X_i=(W_i,Y_i)\), ordered as \((0,0),(0,1),(1,0),(1,1)\), have this joint law. Their one-subject probabilities are
\[
\bigl(q(1-\mu_0),q\mu_0,p(1-\mu_1),p\mu_1\bigr).
\]
\item \textbf{(Comparator noise.)} The pairs \((X_i,L_i)\) are jointly independent across subjects, each \(L_i\) is independent of \(X_i\), and
\[
L_i\sim\operatorname{Lap}(0,\Delta_A/\varepsilon),
\qquad
\Delta_A=\frac1p+\frac1q.
\]
\item \textbf{(Privacy and critical value.)} We have \(\varepsilon>0\), \(0<\gamma<1\), and
\[
z_{1-\gamma/2}
=\inf\{z\in\mathbb R:\Phi(z)\ge1-\gamma/2\}>0,
\]
where \(\Phi\) is the standard-normal distribution function.
\item \textbf{(Strong-saddle selection.)} Fix any measurable strong-saddle selector
\(\vartheta\mapsto(\alpha_\vartheta,t_\vartheta,J_\vartheta)\)
as in \cref{obj:def:pilot-estimator}. In particular, at every interior \(\vartheta\), its selected weights maximize the profiled objective at its selected minimizing coordinate:
\[
F_\vartheta(\alpha,t_\vartheta)
\le F_\vartheta(\alpha_\vartheta,t_\vartheta)
\qquad\text{for every }\alpha\in\mathcal A_\varepsilon.
\]
\end{itemize}

Define the comparator releases, their sample mean, and their empirical variance by
\[
\widetilde A_i
=\frac{W_iY_i}{p}-\frac{(1-W_i)Y_i}{q}+L_i,
\qquad
\widehat\tau_{\mathrm{OA}}
=\frac1n\sum_{i=1}^n\widetilde A_i,
\qquad
\widehat V_{\mathrm{OA}}
=\frac1n\sum_{i=1}^n\widetilde A_i^2
-\widehat\tau_{\mathrm{OA}}^2,
\]
and set
\[
V_{\mathrm{OA}}(\theta,p,\varepsilon)
=\frac{\mu_1}{p}+\frac{\mu_0}{q}-\tau^2
+\frac{2\Delta_A^2}{\varepsilon^2}.
\]
The specified public experiment and its private-record representation hold jointly with the following conclusions:
\[
\sqrt n\bigl(\widehat\tau_{\mathrm{OA}}-\tau\bigr)
\rightsquigarrow
N\!\left(0,V_{\mathrm{OA}}(\theta,p,\varepsilon)\right),
\]
\[
\frac{z_{1-\gamma/2}\sqrt{V_{\mathrm{OA}}(\theta,p,\varepsilon)/n}}
     {z_{1-\gamma/2}\sqrt{V^*(\theta,p,\varepsilon)/n}}
\longrightarrow
\sqrt{\frac{V_{\mathrm{OA}}(\theta,p,\varepsilon)}
           {V^*(\theta,p,\varepsilon)}}.
\]
Moreover, for every \(\delta>0\),
\[
\Pr\!\left(
\left|\widehat V_{\mathrm{OA}}
      -V_{\mathrm{OA}}(\theta,p,\varepsilon)\right|>\delta
\right)\longrightarrow0.
\]

There exists a pilot schedule satisfying
\cref{obj:ass:pilot-diverges,obj:ass:pilot-sublinear}, witnessed by
\[
m_n=\lfloor\sqrt n\rfloor.
\]
For this schedule, use the selected pilot procedure in
\cref{obj:def:pilot-estimator} and its variance statistic \(\widehat V^*\).
For every sequence of couplings whose first marginal is the trial probability law and whose second marginal is the transcript law of that pilot procedure at \((\theta,p)\), and for every \(\delta>0\), the coupled probabilities satisfy
\[
\Pr\!\left(
\left|
\sqrt{\frac{\widehat V_{\mathrm{OA}}}{\widehat V^*}}
-\sqrt{\frac{V_{\mathrm{OA}}(\theta,p,\varepsilon)}
             {V^*(\theta,p,\varepsilon)}}
\right|>\delta
\right)\longrightarrow0.
\]

At \((p,\mu_0,\mu_1,\varepsilon)=(1/2,1/2,1/2,1)\),
\[
V_{\mathrm{OA}}=34,
\qquad
V^*=\left(\frac{\cosh(1/2)}{\sinh(1/2)}\right)^2
< V_{\mathrm{OA}}.
\]
\end{theoremv}

\begin{lemmav}{lem:staircase-refinement}[Staircase refinement and score equality]
Let \(p\) be the treatment-assignment probability, let \(\theta=(\mu_0,\mu_1)^\top\) be the arm-mean parameter, and let \(\varepsilon\) be the privacy level. Suppose that:
\begin{itemize}
\item \textbf{(Interior assignment.)} \(p\) satisfies \cref{obj:ass:interior-assignment}.
\item \textbf{(Interior means.)} \(\theta\) satisfies \cref{obj:ass:interior-means}.
\item \textbf{(Fixed privacy.)} The privacy sequence and \(\varepsilon\) satisfy \cref{obj:ass:fixed-privacy}.
\item \textbf{(Stationary privacy.)} \(Q\) is a probability kernel from the four-record alphabet \(\mathcal X\) to an arbitrary measurable space \(Z\), and, for every measurable \(A\subseteq Z\) and every \(x,x'\in\mathcal X\),
\[
Q(A\mid x)\le e^\varepsilon Q(A\mid x').
\]
\end{itemize}
Index the fourteen nonempty proper subsets of \(\mathcal X\) by \(\mathcal S\). The feasible staircase weights and their channel are given by
\[
\mathcal A_\varepsilon
=
\left\{
\alpha:
\alpha_S\ge0\ \text{for every }S\in\mathcal S,\quad
\sum_{S\in\mathcal S}\alpha_S
\begin{cases}
e^\varepsilon,&x\in S,\\
1,&x\notin S
\end{cases}
=1\ \text{for every }x\in\mathcal X
\right\},
\qquad
Q_\alpha(S\mid x)=\alpha_S
\begin{cases}
e^\varepsilon,&x\in S,\\
1,&x\notin S.
\end{cases}
\]
Writing \(\pi_{\theta,j}=P_\theta(X=x_j)\) for the record probabilities, define the pattern masses, gradients, profiling directions, and projected scores by
\[
h_S(\theta)
=
\sum_{j=1}^{4}\pi_{\theta,j}
\begin{cases}
e^\varepsilon,&x_j\in S,\\
1,&x_j\notin S,
\end{cases}
\qquad
g_S=\nabla_\theta h_S(\theta),
\qquad
v(t)=\begin{pmatrix}t\\t+1\end{pmatrix},
\qquad
\rho_S(\theta,t)=\frac{g_S^\top v(t)}{h_S(\theta)},
\]
with \(p\) and \(\varepsilon\) fixed in the gradient. Let \(I_\theta(Q)\) denote the output Fisher-information matrix, obtained by integrating the outer product of the output score against the output law, and set
\[
I_\theta(\alpha)
=
\sum_{S\in\mathcal S}
\frac{\alpha_S g_Sg_S^\top}{h_S(\theta)}.
\]
Then there exist \(\alpha\in\mathcal A_\varepsilon\) and a probability kernel \(K\) from \(\mathcal S\) to \(Z\) such that
\[
Q=K\circ Q_\alpha,
\qquad
I_\theta(\alpha)-I_\theta(Q)\succeq0.
\]
For every \(t\in\mathbb R\), if
\[
v(t)^\top I_\theta(\alpha)v(t)
=
v(t)^\top I_\theta(Q)v(t),
\]
then every pair \(S,T\in\mathcal S\) with \(\alpha_S\ne0\), \(\alpha_T\ne0\), and conditional output measures \(K(S,\cdot)\) and \(K(T,\cdot)\) that are not mutually singular satisfies
\[
\rho_S(\theta,t)=\rho_T(\theta,t).
\]
\end{lemmav}

\begin{definitionv}{synth_7}[Smooth prior density \(w_R\)]
For \(R>0\), define the Lebesgue density \(w_R\) by
\[
w_R(a)=
\begin{cases}
\displaystyle\frac{15}{16R}\left(1-\frac{a^2}{R^2}\right)^2,
& |a|<R,\\[6pt]
0,& |a|\ge R.
\end{cases}
\]
This density is nonnegative, is positive exactly on \((-R,R)\), has closed support \([-R,R]\), and satisfies
\[
\int_{\mathbb R}w_R(a)\,da
=\int_{-R}^{R}w_R(a)\,da=1.
\]
It is continuously differentiable on \(\mathbb R\), with
\[
w_R'(a)=
\begin{cases}
\displaystyle-\frac{15a}{4R^3}\left(1-\frac{a^2}{R^2}\right),
& |a|<R,\\[6pt]
0,& |a|\ge R,
\end{cases}
\]
so both the density and its derivative vanish at the endpoints. Its squared score is integrable under the prior, with
\[
\int_{-R}^{R}
\frac{\{w_R'(a)\}^2}{w_R(a)}\,da=\frac{10}{R^2},
\]
where the integrand is taken to be zero at points where \(w_R(a)=0\).
\end{definitionv}

\begin{definitionv}{def:sequential-vt-handle}[Sequential van Trees handle]
The sequential van Trees handle \(\mathfrak H_{\mathrm{VT}}\) is defined for parameters \(\theta,p,\varepsilon,R\) satisfying:
\begin{itemize}
\item \textbf{(Interiority.)} \cref{obj:ass:interior-assignment,obj:ass:interior-means}.
\item \textbf{(Privacy.)} \(\varepsilon>0\).
\item \textbf{(Prior radius.)} \(R>0\).
\end{itemize}
It is the direction--prior pair
\[
\mathfrak H_{\mathrm{VT}}=(v(t^*),w_R),
\qquad
v(t^*)=(t^*,t^*+1)^{\top},
\]
where \(t^*\) is the selected minimizing coordinate supplied by \cref{obj:thm:finite-oracle}, satisfying
\[
\max_{\alpha\in\mathcal A_\varepsilon}F_\theta(\alpha,t^*)
=J^*(\theta,p,\varepsilon),
\]
with the feasible weights, profiled objective, and optimized value as in \cref{obj:def:staircase-saddle}. The prior density is
\[
w_R(u)=
\begin{cases}
\dfrac{15}{16R}\left(1-\left(\dfrac{u}{R}\right)^2\right)^2,
& |u|<R,\\
0,& |u|\ge R.
\end{cases}
\]
\end{definitionv}
